% !TEX root = ./2nde.tex
\begin{exo}
	Un professeur a relevé puis résumé dans un tableau les notes de ses élèves lors d'un devoir sur $10$.

	\begin{center}
		\renewcommand{\arraystretch}{1.4}
		\begin{tabular}{|l|*{11}{c|}}
			\hline
			\textbf{Note}     & $0$ & $1$ & $2$ & $3$ & $4$ & $5$ & $6$ & $7$ & $8$ & $9$ & $10$ \\ \hline
			\textbf{Effectif} & $1$ & $0$ & $3$ & $4$ & $2$ & $6$ & $7$ & $7$ & $2$ & $2$ & $0$  \\ \hline
		\end{tabular}
	\end{center}

	\begin{enumerate}
		\item Donner le tableau des effectifs cumulés croissants.
		\item Où peut-on lire l'effectif total ? Combien vaut-il ?
		\item Quelle est la note médiane $M_e$ ? Interpréter ce résultat.
		\item Déterminer le premier quartile $Q_1$, puis interpréter le résultat.
		\item Déterminer le troisième quartile $Q_3$.
		\item Déterminer l'intervalle interquartile puis interpréter le résultat.
	\end{enumerate}
\end{exo}
